%!TEX root = main.tex
\setcounter{chapter}{5}
\setcounter{section}{6}
\setcounter{theorem}{0}
\setcounter{equation}{0}

\lectureheader{161}{Calculus I}{Substitution and area between curves}{\textit{Thomas' Calculus} \thesection}

%\begin{remark}
%Units matter, and hence so do differentials.
%\end{remark}
%
%\begin{example}
%When compressed $x$ meters beyond its natural length, a certain spring exerts a force of 
%\begin{equation*}
%F=-20x
%\end{equation*}
%newtons along the $x$-axis.
%An $80$-gram weight is attached to the end of a spring and the spring is compressed and then released.
%It is observed that the position of the weight along the $x$-axis is
%\begin{equation*}
%x = -\frac{1}{5}\cos(t/2)\quad (0\le t\le \pi)
%\end{equation*}
%where $t$ is time (measured in seconds).
%Note that this means that the force exerted by the spring at time $t$ seconds is
%\begin{equation*}
%F=4\cot(t/2)\quad (0\le t\le \pi)
%\end{equation*}
%newtons.
%\begin{enumerate}
%\item Why is it wrong to say that the work done by the spring on the weight over the first $\pi$ seconds is $\DS\int_0^\pi F\dee t = \int_0^\pi 4\cos(t/2)\dee t$?
%\begin{flushright}
%\includegraphics[width=3in]{force_v_time}
%\end{flushright}
%\end{enumerate}
%\end{example}
%
%\newpage

\begin{theorem}[Substitution rule]
If $u=g(x)$ is a differentiable function whose range is in an interval $I$, and $f$ is continuous on $I$, then
\begin{equation*}
\int_a^b f(g(x))g'(x)\dee x = \int_{g(a)}^{g(b)} f(u)\dee u.
\end{equation*}
\end{theorem}

\begin{example}
Evaluate $\DS\int_{-1}^13x^2\sqrt{x^3+1}\dee x$.
\end{example}

\newpage

\begin{example}
Evaluate $\DS\int_{\pi/4}^{\pi/2}\cot\theta\csc^2\theta\dee \theta$.
\end{example}

\vfill

\begin{example}
Evaluate $\DS\int_{0}^{\pi/2}\frac{2\sin x\cos x}{(1+\sin^2 x)^3}\dee x$.
\end{example}

\vfill

\newpage

\begin{theorem}
Suppose that $f$ is continuous on $[-a, a]$.
\begin{enumerate}
\item If $f$ is even, then $\DS\int_{-a}^a f(x)\dee x = 2\int_0^af(x)\dee x$.
\item If $f$ is odd, then $\DS\int_{-a}^a f(x)\dee x = 0$.
\end{enumerate}
\end{theorem}

\begin{example}
Evaluate $\DS\int_{-2}^2(x^4-4x^2+6)\dee x$.
\end{example}

\newpage

\begin{definition}
If $f$ and $g$ are continuous on $[a,b]$, then the \textbf{area} $A$ bounded by the curves $y=f(x)$ and $y=g(x)$  and the lines $x=a$ and $x=b$ is 
\begin{equation*}
A = \int_a^b\Big|f(x)-g(x)\Big|\dee x.
\end{equation*}
\end{definition}
\begin{remark}
In practice, it is usually a good idea to make a rough sketch of the area before trying to set up an integral to compute its area.
\end{remark}

\begin{example}
Find the area of the region enclosed by the parabola $y=2-x^2$ and the line $y=-x$.
\end{example}

\newpage

\begin{example}
Find the area of the region in the first quadrant that is bounded above by $y=\sqrt x$ and below by the $x$-axis and the line $y=x-2$.
\end{example}

\newpage

\begin{example}
Repeat the previous exercise using integration with respect to $y$.
\end{example}

\newpage

\begin{example}
Find the area of the region bounded below by the line $y=2-x$ and above by the curve $y=\sqrt{2x-x^2}$.
\end{example}

